% frontmatter/palabras/verano.tex -- instrucciones para el adulto del
% cuaderno de verano de primeras palabras (ediciones/palabras-verano.tex,
% ver notes/07-ediciones.md): van en vez de
% frontmatter/palabras/como-usar.tex, que habla del año entero. Detrás,
% en vez del mapa del curso, el mapa del verano (\mapaEdicion,
% preamble.tex).
\section*{Cómo usar este cuaderno de verano}

Este cuaderno es el verano de \emph{Primeras palabras}, el primer
cuaderno de \emph{Aprendo a leer}: sus últimas 13 semanas, en un
cuaderno aparte. Es el verano en el que se pasa de leer palabras a leer
frases: cada día, una frase corta, de dos a cuatro palabras
(\emph{Toby espera.}, \emph{Bigotes duerme.}). Está pensado para una
niña o un niño que ya lee palabras de varias sílabas -- también con
sílabas trabadas (\emph{pla, tre}), con \emph{ch, ll, rr, qu, gu} y con
diptongos (\emph{ue, ie}) -- y que empieza a juntarlas en frases. Al
terminarlo, estará justo donde empieza el cuaderno de frases.

Hay una página para cada día de lunes a viernes, 65 en total. Cinco o
diez minutos al día son suficientes. Cada página sigue siempre la misma
estructura:

\begin{itemize}[leftmargin=6mm]
\item \textbf{Fecha}, \textbf{Leído} y \textbf{Notas}: para el adulto.
  Marcar cada día si ha leído sola, con ayuda o con dificultad no
  cuesta nada, y al final del verano es un registro de lo que ha
  avanzado.
\item La \textbf{frase del día}, en el recuadro verde y en letra
  grande: se lee señalando cada palabra con el dedo.
\item Una \textbf{actividad}, que casi siempre termina dibujando. Las
  instrucciones (en letra pequeña y gris) \textbf{las lee un adulto} en
  voz alta: lo único que tiene que leer la niña o el niño es lo que
  está en letra grande.
\end{itemize}

\subsection*{Las actividades}

{\small
\begin{itemize}[leftmargin=6mm, itemsep=1pt, topsep=2pt]
\item \textbf{\lblDibuja}: dibuja lo que acabas de leer. Es la manera
  más sencilla de comprobar que la frase se ha entendido, no solo
  descifrado.
\item \textbf{\lblCompleta}: unas pocas líneas sin forma que hay que
  convertir en un dibujo.
\item \textbf{\lblTraza}: repasar con el dedo o con un lápiz una letra,
  mayúscula y minúscula, y después escribirla sola en la línea.
  \textbf{\lblEscribe}: lo mismo con una palabra entera, por dentro de
  sus letras huecas.
\item \textbf{\lblEncuentra}: rodear una palabra cada vez que aparece
  entre otras que se le parecen mucho. Obliga a mirar todas las letras,
  no solo la primera.
\item \textbf{\lblAdivina}: el adulto lee una adivinanza, y la niña o
  el niño dibuja la respuesta y escribe su nombre debajo.
\item \textbf{\lblUne}: unir con una línea cada animal con su sonido,
  cada palabra con su contraria o con la que rima.
\item \textbf{\lblSiNo}: leer frases muy cortas y decidir si tienen
  sentido (\emph{Toby ladra} -- sí; \emph{Toby vuela} -- no).
\item \textbf{\lblRelee} y \textbf{\lblRepasa} (los viernes): las frases
  de toda la semana juntas, para leerlas otra vez y marcar las que ya
  salen solas. Cada diez páginas hay un pequeño cartel de celebración.
\end{itemize}}

\subsection*{Cómo ayudar}

{\small
\begin{itemize}[leftmargin=6mm, itemsep=1pt, topsep=2pt]
\item Si se atasca en una palabra, déjale unos segundos y, si no sale,
  léela tú señalando las sílabas con el dedo -- y que la repita. Mejor
  eso que convertir la página en una lucha.
\item No hace falta corregir cada error. Si ha leído \emph{pato} en vez
  de \emph{pata}, basta con preguntar ``¿seguro? mira el final''.
\item La actividad no tiene que terminarse el mismo día. Lo importante
  es la lectura; el dibujo puede esperar a la tarde.
\end{itemize}}

Las páginas cuentan el verano de \textbf{Lucía}, \textbf{Dani} y su
perro \textbf{Toby} en el pueblo de la abuela Rosa, con las mismas
semanas que el cuaderno de verano de frases: dos hermanos pueden leer
cada uno el suyo sobre lo mismo. En la página siguiente están todas, con
un sol para colorear cada día que se lee; al final del cuaderno hay una
clave de respuestas y un diploma.
\newpage
